\begin{table}[htbp]
\caption{Treatment-Control Differences in Selected Covariates and Sample Indicators: Kenya Year 1}
\label{tab:ke_balance}
\begin{center}
\begin{small}
\begin{threeparttable}
\begin{tabular}{l*{5}{r}}
\toprule
 & \Longstack{Treatment\\ mean} & \Longstack{Control\\ mean} & Difference & \Longstack{Standard\\ error} & $p$-value \\
\midrule
Baseline score & 39.904 & 39.207 & 0.742 & (0.657) & 0.259 \\
Upper group & 0.652 & 0.614 & 0.040 & (0.033) & 0.218 \\
Has endline & 0.771 & 0.686 & 0.075 & (0.045) & 0.095 \\
Class size & 20.935 & 24.307 & -3.597 & (1.770) & 0.042 \\
\bottomrule
\end{tabular}
\begin{tablenotes}
\footnotesize \item Notes: The unit of observation is a student in the Kenya Year 1 analytic sample. The table reports treatment and control means together with the treatment-control difference from a student-level OLS regression that includes constituency-strata fixed effects. Standard errors are cluster-robust at the school-grade-group level. `Upper group' is the share of students whose baseline score lies above the Kenya Year 1 grade-specific reading-group cutoff, `Has endline' is the share with a matched endline outcome, and `Class size' is realized reading-class size. The table is descriptive and is used to assess treatment-control comparability on key observable margins.
\end{tablenotes}
\end{threeparttable}
\end{small}
\end{center}
\end{table}